 \documentclass[a4paper,twoside,11pt]{mwrep}
\usepackage[utf8]{inputenc}
\usepackage{polski}
\usepackage{graphicx}
\usepackage{xcolor}
%%%%%%%%%%%%%%%%%%%%%%%%%%%%
\usepackage{hyperref}
\hypersetup{
    colorlinks=true,
    linkcolor=black!30!blue,
    filecolor=black!30!magenta, 
    citecolor=black!30!green,      
    urlcolor=black!30!cyan
}
%%%%%%%%%%%%%%%%%%%%%%%%%%%
\usepackage{amsmath}
\usepackage{amsthm}
\usepackage{amssymb}
\usepackage{tikz}
\usetikzlibrary{positioning, arrows.meta}
\usepackage{booktabs}
\usepackage{longtable}
\usepackage{listings}
\usepackage{caption}
%%%%%%%%%%%%%%%%%%%%%%%%%%%%
%%%%%%%%%%%%%%%%%%%%%%%%%%%%
\usepackage{geometry}
 \geometry{
    paper=a4paper,
    inner=30mm,
    outer=20mm,
    nohead,
    top=20mm,
    bottom=25mm,
    }
%%%%%%%%%%%%%%%%%%%%%%%%%%%
\newcommand\blankpage{%
    \null
    \thispagestyle{empty}%
    \addtocounter{page}{-1}%
    \newpage}
%%%%%%%%%%%%%%%%%%%%%%%%%%
\theoremstyle{plain} \newtheorem{thm}{Twierdzenie}
\theoremstyle{plain} \newtheorem{lem}{Lemat}
\theoremstyle{definition} \newtheorem{defn}{Definicja}
\theoremstyle{remark} \newtheorem*{rem}{Wniosek}

\lstset{
    basicstyle=\ttfamily\footnotesize,
    breaklines=true,
    frame=single,
    columns=fullflexible,
    captionpos=b
}
%%%%%%%%%%%%%%%%%%%%%%%%%%%%%%%%%%%%%%%%%%%%%%%%%%%%%%%%%%%%%%%%
\newcommand{\autor}{\fontfamily{qhv}\fontshape{n}\selectfont\LARGE\bfseries
Samuel Rus}
\newcommand{\album}{\fontfamily{qhv}\fontshape{n}\selectfont
\ldots\ldots}
\newcommand{\pltitle}{\fontfamily{qhv}\fontshape{n}\selectfont\LARGE\bfseries

Wykorzystanie wybranych metod uczenia maszynowego w predykcji efektywności sportowej zawodników Premier League.}
\newcommand{\engtitle}{\fontfamily{qhv}\fontshape{n}\selectfont\LARGE\bfseries

Predicting Premier League Player Performance Using Selected Machine Learning Methods}
\newcommand{\level}{\fontfamily{qhv}\fontshape{n}\selectfont\LARGE
inżynierska}
\newcommand{\kierunek}{\fontfamily{qhv}\fontshape{n}\selectfont\LARGE
Matematyka stosowana}
\newcommand{\promotor}{\fontfamily{qhv}\fontshape{n}\selectfont\large\bfseries
dr Adam Sebastian Marszałek}
\newcommand{\spec}{\fontfamily{qhv}\fontshape{n}\selectfont\LARGE\bfseries
 Analiza danych}
\newcommand{\rok}{\fontfamily{qhv}\fontshape{n}\selectfont
202\ldots}

%%%%%%%%%%%%%%%%%%%%%%%%%%%%%%%%%%%%%%%%%%%%%%%%%%%%%%%%%%%%%%%%
%%%%%%%%%%%%%%%%%%%%%%%%%%%%%%%%%%%%%%%%%%%%%%%%%%%%%%%%%%%%%%%%
\begin{document}
 \begin{titlepage}
 \thispagestyle{empty}
 \vspace*{3ex}
\fontfamily{qhv}\fontshape{n}\selectfont
\hspace{-1.3em}\parbox{0.15\textwidth}{\includegraphics[height=0.15\textwidth]{logoPK.png}}
\parbox[c]{0.7\textwidth}
{\centering
\LARGE
{\bfseries Politechnika Krakowska}\\
{\bfseries im. Tadeusza Kościuszki}\\

\smallskip

\LARGE Wydział Informatyki i Matematyki}\hspace{0.3em}
 \parbox{0.2\textwidth}{\includegraphics[width=0.15\textwidth]{logoIT.png}}
\parbox[c]{0.15\textwidth}

\vspace{2ex}

 \vspace{0.13\textheight}

 \center{\autor}
 \smallskip
 \center{\large numer albumu: 151814 }

\vspace*{3ex}

 \center{\pltitle}
 
 \smallskip

 \center{\engtitle}

 \bigskip

 \center{\LARGE\bfseries praca \level \\[1ex]
\fontfamily{qhv}\fontshape{n}\selectfont\LARGE\bfseries na kierunku
 \kierunek\\[1ex]
 specjalność \spec} 

\vspace{0.18\textheight}


\hspace{0.4\textwidth} \parbox{0.5\textwidth}{Praca przygotowana pod
kierunkiem:
\\
\bfseries \promotor}


\vfill

 \center{Kraków, 2026/27}
 \end{titlepage}
\newpage
\blankpage
%%%%%%%%%%%%%%%%%%%%%%%%%%%%%%%%%%%%%%%%%%%%%%%%%%%%%%%%%%%%%%%%%%%%%%%%%%%%%%
\pagenumbering{roman}

\thispagestyle{empty}
\chapter*{Streszczenie}
\addcontentsline{toc}{chapter}{Streszczenie}
\noindent
Praca dotyczy predykcji punktacji zawodników w grze Fantasy Premier League przy użyciu metod uczenia maszynowego. Zaproponowany pełny proces
przetwarzania danych, łączący dane historyczne kolejka po kolejce oraz kursy bukmacherskie przeliczone na prawdopodobieństwa. Po eksploracyjnej analizie
danych oraz inżynierii cech
respektującej rygor braku wycieku informacji z przyszłości,
porównano trzy rodziny modeli regresyjnych: uogólnione modele liniowe, modele zespołowe oparte na drzewach decyzyjnych oraz sieci neuronowe. Zaprojektowano autorską funkcję straty oraz
metryki ewaluacyjne segmentowane względem ceny rynkowej zawodnika, odzwierciedlające asymetryczną
użyteczność decyzyjną błędu predykcji z perspektywy użytkownika gry.
Rozważono również rozszerzenie architektury o dodatkową warstwę (dalej nazywaną meta modelem), łączącą predykcje modeli bazowych z poszanowaniem chronologicznego podziału danych, wraz z oceną wpływu takiego podejścia na jakość i wiarygodność predykcji.

\clearpage
\thispagestyle{empty}
\chapter*{Abstract}
\addcontentsline{toc}{chapter}{Abstract}
\noindent
This thesis addresses the prediction of Fantasy Premier League (FPL)
player point scores using machine learning methods. A complete data
pipeline is proposed, combining historical gameweek-level data (the
\emph{vaastav/Fantasy-Premier-League} repository) and bookmaker odds
(\emph{football-data.co.uk}) converted into probabilities. Following exploratory data
analysis (target distribution analysis, Spearman rank correlations,
PCA, k-means clustering with Kruskal--Wallis testing) and feature
engineering designed to strictly avoid data leakage (pre-game features,
exponential moving averages, lagged features, rolling standard
deviations), three families of regression models are compared:
generalized linear models (GLM), tree-based ensemble models (Random
Forest, XGBoost, CatBoost), and neural networks (MLP, a 1D convolutional
network, and an LSTM network). A custom loss function (\texttt{FPLLoss})
and price-segmented evaluation metrics (budget / mid / premium players)
are proposed, reflecting the asymmetric decision-relevant cost of
prediction errors from the perspective of an FPL player. An extension of
the architecture with an additional stacking layer (a meta-model
combining base-model predictions under a strict chronological data
split) is also considered, together with an assessment of its impact on
predictive quality and reliability.

\medskip
\noindent\textbf{Keywords:} Fantasy Premier League, machine learning,
regression, generalized linear models, XGBoost, CatBoost, neural
networks, LSTM, loss function, stacking, sports analytics.

%%%%%%%%%%%%%%%%%%%%%%%%%%%%%%%%%%%%%%%%%%%%%%%%%%%%%%%%%%%%%%%%%%%%%%%%%%%%%%
\tableofcontents
%%%%%%%%%%%%%%%%%%%%%%%%%%%%%%%%%%%%%%%%%%%%%%%%%%%%%%%%%%%%%%%%%%%%%%%%%%%%%%
\pagenumbering{arabic}
\setcounter{page}{1}

%======================================================================
\chapter*{Wstęp}
\addcontentsline{toc}{chapter}{Wstęp}

Fantasy Premier League jest jedną z najpopularniejszych gier sportowych na świecie, w której uczestnicy zarządzają wirtualną
drużyną piłkarzy angielskiej Premier League, otrzymując punkty na podstawie
rzeczywistych występów zawodników w danej kolejce ligowej. Gra angażuje kilkanaście milionów użytkowników w każdym sezonie, co sprawia że problem predykcji liczby punktów zdobywanych przez zawodnika w kolejnej
kolejce stał się interesującym zagadnieniem badawczym,
łączącym elementy analizy szeregów czasowych, regresji oraz uczenia maszynowego
stosowanego do danych sportowych.

Predykcja punktacji FPL charakteryzuje się kilkoma cechami utrudniającymi
zastosowanie standardowych metod regresji:
\begin{itemize}
    \item Rozkład zmiennej celu jest silnie asymetryczny i skoncentrowany
    wokół zera. Znaczna część zawodników w danej kolejce nie wychodzi na
    boisko lub gra niewielką liczbę minut.
    \item Wariancja punktacji jest silnie zróżnicowana w zależności od
    pozycji oraz wartości rynkowej zawodnika: gracze o wysokiej cenie mają
    zarówno wyższą oczekiwaną liczbę punktów, jak i większą wariancję
    wyników, co czyni błąd predykcji dla tej grupy szczególnie kosztownym
    z perspektywy użytkownika podejmującego decyzje transferowe.
    \item Dane mają naturalną strukturę czasową, co wymusza odpowiednie
    podejście do konstrukcji cech opóźnionych oraz podziału zbioru na
    część treningową, walidacyjną i testową w sposób chronologiczny,
    zapobiegający wyciekowi danych z przyszłości.
\end{itemize}


Celem pracy jest zaprojektowanie, implementacja oraz porównanie kilku
rodzin metod uczenia maszynowego: od uogólnionych modeli liniowych, przez modele oparte na drzewach decyzyjnych, po sieci neuronowe w zadaniu predykcji liczby punktów FPL
zdobywanych przez zawodnika w pojedynczej kolejce ligowej, na podstawie
danych historycznych, statystyk oczekiwanych, kursów
bukmacherskich przeliczonych na prawdopodobieństwa. Wszystkie dane pochodzą z otwartych źródeł i są ujednolicane do spójnego formatu w procesie integracji. Istotnym elementem pracy jest zaprojektowanie funkcji straty
oraz metryk ewaluacyjnych uwzględniających segmentację zawodników według
ceny rynkowej, tak aby jakość predykcji była oceniana w sposób zbliżony do
rzeczywistej użyteczności decyzyjnej dla gracza FPL.

Praca obejmuje pełny proces przetwarzania danych: od automatycznego
pozyskania danych historycznych z publicznych repozytoriów i serwisów, kursy bukmacherskie, przez eksploracyjną
analizę danych i inżynierię cech, po budowę, trenowanie i ewaluację modeli
predykcyjnych wraz ze śledzeniem eksperymentów. Dodatkową częscią pracy, będzie wybór składu na daną kolejke Premier League.


%======================================================================
\chapter{Charakterystyka problemu}\label{ch:problem}
\section{Zasady gry Fantasy Premier League}

W grze FPL uczestnik tworzy drużynę piłkarzy Premier League, mając do dyspozycji 100 milionów funtów wirtualnego budżetu na początku sezonu. Wybór obejmuje bramkarzy, obrońców, pomocników oraz napastników. W każdej kolejce ligowej zawodnicy otrzymują punkty na podstawie zdarzeń meczowych: goli, asyst, czystych kont, kartek, interwencji bramkarskich czy statystyk wpływających na system bonusowy. Kluczową cechą systemu punktowego jest jego zależność od nominalnej pozycji zawodnika - bramkarze i obrońcy są znacznie wyżej nagradzani za czyste konto lub strzelonego gola niż gracze formacji ofensywnych, co uzasadnia podział analizy według pozycji, stosowaną w dalszych rozdziałach pracy.

Szczegółowe zestawienie zasad punktacji z uwzględnieniem podziału na pozycje boiskowe przedstawiono w Tabeli~\ref{tab:fpl_punktacja}.

\begin{table}[!htbp]

\renewcommand{\arraystretch}{0.92} 
\resizebox{\linewidth}{!}{%
\begin{tabular}{lcccc}
\hline
\textbf{Zdarzenie meczowe} & \textbf{Bramkarz} & \textbf{Obrońca} & \textbf{Pomocnik} & \textbf{Napastnik   } \\
\hline
Gra do 60 minut & 1 & 1 & 1 & 1 \\
Gra co najmniej 60 minut & 2 & 2 & 2 & 2 \\
Gol & 6 & 6 & 5 & 4 \\
Asysta & 3 & 3 & 3 & 3 \\
Czyste konto (min. 60 min gry) & 4 & 4 & 1 & 0 \\
Każde 3 obronione strzały & 1 & -- & -- & -- \\
Obrona rzutu karnego & 5 & -- & -- & -- \\
Niewykorzystany rzut karny & -2 & -2 & -2 & -2 \\
Każde 2 stracone bramki & -1 & -1 & 0 & 0 \\
Żółta kartka & -1 & -1 & -1 & -1 \\
Czerwona kartka & -3 & -3 & -3 & -3 \\
Gol samobójczy & -2 & -2 & -2 & -2 \\
Punkty bonusowe (BPS)* & 1--3 & 1--3 & 1--3 & 1--3 \\
\hline
\multicolumn{5}{l}{\footnotesize{* Punkty bonusowe przyznawane 3 najlepszym graczom wg wskaźnika BPS.}} \\

\end{tabular}%
}
\end{table}
\section{Sformułowanie zadania predykcyjnego}

Zadanie zdefiniowano jako regresję jednowymiarową: dla zawodnika $i$ w
kolejce $t$ sezonu $s$ celem jest predykcja wartości
$y_{i,t,s} = \texttt{total\_points}_{i,t,s} \in \mathbb{Z}_{\geq 0}$
na podstawie wektora cech $\mathbf{x}_{i,t,s}$ zbudowanego wyłącznie z
informacji dostępnych \emph{przed} rozpoczęciem kolejki $t$ (tzw. cechy
\emph{pre-game}), tj. bez wykorzystania statystyk z samego meczu, którego
wynik jest przedmiotem predykcji. Ograniczenie to (opisane szczegółowo w
rozdziale~\ref{ch:features}) jest krytyczne dla poprawności metodologicznej
-- jego naruszenie prowadziłoby do wycieku danych i sztucznie zawyżonej
jakości modelu w warunkach walidacji, niemożliwej do odtworzenia w
rzeczywistym użyciu prognozy przed kolejką.

Dodatkowym utrudnieniem jest silna heterogeniczność zadania względem ceny
rynkowej zawodnika (\texttt{value}), będącej w przybliżeniu \emph{proxy} dla
jego jakości i regularności gry. W pracy przyjęto podział zawodników na
trzy segmenty cenowe -- budżetowy, średni i premium -- na podstawie
kwantyli rozkładu ceny (odpowiednio 50. i 95,5. percentyl), co pozwala
raportować jakość predykcji oddzielnie dla każdej grupy (rozdział
\ref{ch:loss}).

%======================================================================
\chapter{Źródła i pozyskiwanie danych}\label{ch:dane}

\section{Przegląd źródeł danych}

Dane wykorzystane w pracy pochodzą z trzech niezależnych źródeł, łączonych
następnie w jeden spójny zbiór przy użyciu procedur ujednolicania formatów:

\begin{itemize}
    \item \textbf{vaastav/Fantasy-Premier-League} -- publiczne
    repozytorium GitHub agregujące historyczne dane kolejka po kolejce oraz metadane zawodników (\texttt{players\_raw.csv})
    dla sezonów Premier League od 2022/23 do bieżącego sezonu 2025/26;
    \item \textbf{football-data.co.uk} -- serwis udostępniający historyczne
    kursy bukmacherskie (rynek 1X2 oraz powyżej/poniżej 2,5 gola) dla
    poszczególnych spotkań ligowych;
\end{itemize}

\section{Pozyskiwanie danych kolejka-po-kolejce (\texttt{FetchVaastav})}

Klasa \texttt{FetchVaastav} pobiera dane dla każdego sezonu i każdej z 38
kolejek osobno, łącząc je następnie z metadanymi tożsamości zawodników
(\texttt{id}, imię i nazwisko, \texttt{code}) po kluczu
\texttt{(element, season)}. W procesie czyszczenia danych
(\texttt{\_merge\_and\_clean}) wykonywane są następujące operacje:
\begin{itemize}
    \item normalizacja ceny zawodnika (\texttt{value}) do jednostek
    właściwych grze (dzielenie przez 10);
    \item normalizacja popularności wyboru (\texttt{selected}) do rozkładu
    prawdopodobieństwa w obrębie każdej pary (sezon, kolejka), poprzez
    podział przez sumę w grupie;
    \item parsowanie i ujednolicenie znacznika czasu meczu
    (\texttt{kickoff\_time});
    \item ujednolicenie etykiety sezonu do formatu czterocyfrowego (np.
    \texttt{"2022-23"} $\rightarrow$ \texttt{2223});
    \item usunięcie kolumn nieistotnych merytorycznie lub wprowadzonych
    zmianami w strukturze danych źródłowych między sezonami (m.in. cechy
    związane z rolą menedżera drużyny, \texttt{mng\_*});
    \item odrzucenie rekordów z brakami w kluczowych kolumnach
    identyfikacyjnych (\texttt{team}, \texttt{season}, \texttt{was\_home},
    \texttt{kickoff\_time}).
\end{itemize}

Dodatkowo zaimplementowano metodę walidacyjną \texttt{run\_tests},
realizującą automatyczny raport jakości danych (\emph{data quality report}):
kompletność liczby kolejek w sezonie (oczekiwane 38), odsetek braków
danych w poszczególnych kolumnach, spójność schematu kolumn między
sezonami, zakres wartości ceny zawodnika oraz wykrywanie duplikatów po
kluczu głównym \texttt{(code, season, gw)}.

\section{Pozyskiwanie kursów bukmacherskich (\texttt{FetchBetting})}

Klasa \texttt{FetchBetting} pobiera dane kursów bukmacherskich (rynek
zwycięstwo gospodarzy / remis / zwycięstwo gości oraz liczba goli
powyżej/poniżej 2,5) i przekształca je w dwóch krokach:

\begin{enumerate}
    \item \textbf{Konwersja kursów na prawdopodobieństwa} -- surowe
    prawdopodobieństwo implikowane obliczane jest jako odwrotność kursu
    ($\hat{p}_k = 1/\text{kurs}_k$), a następnie normalizowane w celu
    usunięcia marży bukmacherskiej (\emph{overround}):
    \begin{equation}
        p_k = \frac{\hat{p}_k}{\sum_{j} \hat{p}_j}, \qquad \sum_k p_k = 1.
    \end{equation}
    \item \textbf{Transformacja z formatu szerokiego do długiego} -- każdy
    mecz (gospodarz--gość) jest rozbijany na dwa rekordy, jeden z
    perspektywy drużyny gospodarzy, drugi z perspektywy drużyny gości, z
    odpowiednio dobranymi prawdopodobieństwami zwycięstwa/remisu/porażki
    oraz opcjonalnym handicapem azjatyckim (\texttt{ahh}), którego znak jest
    odwracany dla drużynygości.
\end{enumerate}


\section{Łączenie źródeł (\texttt{MergeData})}
\begin{figure}[htbp]
    \centering
    \begin{tikzpicture}[
        font=\sffamily\small,
        source/.style={
            draw=blue!40!black, fill=blue!4, thick, rounded corners=8pt,
            align=left, text width=5.8cm, minimum height=5.8cm,
            inner sep=10pt, text=black!90
        },
        merge/.style={
            draw=green!45!black, fill=green!6, thick, rounded corners=8pt,
            align=center, text width=8cm, minimum height=2cm,
            inner sep=12pt, text=black!90
        },
        arrow/.style={
            ->, thick, >=stealth, draw=gray!80!black, rounded corners=8pt
        }
    ]

    \node[source] (vaastav) {
        \textbf{\large Vaastav (repozytorium GitHub)} \\[0.2cm]
        \textbf{Opis:} Publiczne repozytorium agregujące oficjalne statystyki z gry FPL. \\[0.1cm]
        \textbf{Sezony:} 2022/23 -- 2025/26 \\[0.1cm]
        \textbf{Kluczowe cechy:}
        \begin{itemize}
            \setlength\itemsep{0em}
            \item \texttt{id, name, team, code}
            \item \texttt{gw, kickoff\_time, was\_home}
            \item \texttt{value, selected}
        \end{itemize}
    };

    \node[source, right=1cm of vaastav] (fci) {
        \textbf{\large FPL-Core-Insights (FCI)} \\[0.2cm]
        \textbf{Opis:} Zewnętrzna baza dostarczająca metryki analityczne (xG, xA) i statystyki transferowe. \\[0.1cm]
        \textbf{Sezony:} 2025/26 (bieżący) \\[0.1cm]
        \textbf{Kluczowe cechy:}
        \begin{itemize}
            \setlength\itemsep{0em}
            \item \texttt{xP} (oczekiwane punkty)
            \item \texttt{expected\_goals, expected\_assists}
            \item \texttt{transfers\_in, transfers\_out}
        \end{itemize}
    };

    \path (vaastav.south east) -- (fci.south west) coordinate[midway] (mid);

    \node[source, below=1.5cm of mid] (betting) {
        \textbf{\large football-data.co.uk} \\[0.2cm]
        \textbf{Opis:} Serwis gromadzący kursy bukmacherskie z rynków 1X2 oraz linii bramkowych. \\[0.1cm]
        \textbf{Sezony:} 2022/23 -- 2025/26 \\[0.1cm]
        \textbf{Kluczowe cechy:}
        \begin{itemize}
            \setlength\itemsep{0em}
            \item \texttt{prob\_win, prob\_draw, prob\_lose}
            \item \texttt{prob\_over\_2.5, prob\_under\_2.5}
            \item \texttt{ahh} (handicap azjatycki)
        \end{itemize}
    };

    \node[merge, below=1.7cm of betting] (dataset) {
        \textbf{\Large Zintegrowany zbiór danych} \\[0.2cm]
        \normalsize Zapis do formatu Parquet po asynchronicznym \\
        łączeniu (\texttt{merge\_asof}) i złączeniach wewnętrznych
    };

    \draw[arrow] (vaastav.west) -- ++(-0.5cm,0) |- (dataset.west);
    \draw[arrow] (fci.east)     -- ++(0.5cm,0)  |- (dataset.east);
    \draw[arrow] (betting.south) -- (dataset.north);

    \end{tikzpicture}
    \caption{Architektura zasilania modelu FPL. Wszystkie źródła przedstawiono jednolitym formatem.}
    \label{fig:data_sources_clean}
\end{figure}

Finalne połączenie dwóch źródeł danych odbywa się w jednym kroku:
\section{Wstępne przekształcenia}

Przed właściwą analizą wykonano ujednolicenie typów danych (rzutowanie
kolumn liczbowych do \texttt{float32}), zmianę nazw wybranych kolumn
statystyk oczekiwanych na czytelniejsze odpowiedniki (np.
\texttt{expected\_goals} $\rightarrow$ \texttt{xg}) oraz normalizację
statystyk ofensywnych do wartości \emph{na 90 minut} gry
($x_{\text{per90}} = x \cdot 90 / \text{minuty}$), co umożliwia
porównywanie zawodników o różnym czasie gry. Z analizy wykluczono pozycję
\texttt{AM} (napastnik pomocniczy) jako niereprezentowaną w bieżącym
schemacie pozycji gry.

\section{Analiza rozkładu zmiennej celu}

Rozkład zmiennej \texttt{total\_points} scharakteryzowano za pomocą
histogramu, wykresu kwantyl-kwantyl (QQ-plot) względem rozkładu
normalnego oraz miar skośności i kurtozy. Analiza potwierdziła silną
prawostronną skośność rozkładu oraz istotny udział wartości zerowych,
wynikający z zawodników niewystępujących w danej kolejce -- obserwacja ta
uzasadnia wykorzystanie w dalszej części pracy rozkładów właściwych dla
danych zliczeniowych (Poissona, ujemnego dwumianowego, Tweediego) w
modelach GLM (rozdział~\ref{ch:models}) oraz odpowiednio zaprojektowanej
funkcji straty w modelach sieci neuronowych (rozdział~\ref{ch:loss}).

\section{Analiza korelacji}

Ze względu na nieliniowość i obecność wartości odstających zastosowano
współczynnik korelacji rang Spearmana (nie Pearsona), obliczony osobno dla
całego zbioru oraz osobno w podziale na pozycję zawodnika. Dla każdej
pozycji wyznaczono ranking piętnastu (odpowiednio: jedenastu) cech
najsilniej skorelowanych z \texttt{total\_points}, co pozwoliło
zidentyfikować różnice w istotności predyktorów między pozycjami (np.
większe znaczenie cech ofensywnych xG/xA dla pomocników i napastników,
a defensywnych -- dla obrońców i bramkarzy).

\section{Efekt gry u siebie}

Zweryfikowano hipotezę o wpływie czynnika gry u siebie (\texttt{was\_home})
na średnią liczbę punktów, poprzez porównanie średnich wartości
\texttt{total\_points} w obu grupach.

\section{Redukcja wymiarowości i analiza głównych składowych (PCA)}

Do standaryzowanych cech numerycznych zastosowano potok
\texttt{RobustScaler} $\rightarrow$ \texttt{PCA} (skalowanie odporne na
wartości odstające, istotne przy silnie skośnym rozkładzie wielu cech
wydajnościowych). Przeanalizowano wariancję wyjaśnianą przez kolejne
składowe (wykres \emph{scree plot} wraz z wariancją skumulowaną) oraz
wykres ładunków (\emph{loading plot}) dwóch pierwszych składowych głównych,
identyfikujący cechy o największym wkładzie w kierunki głównej zmienności
danych.

\section{Analiza skupień}

Dla każdej pozycji osobno wykonano standaryzację cech (\texttt{RobustScaler})
i grupowanie metodą \emph{k-średnich} (\texttt{KMeans}), z liczbą skupień
$k=3$ (interpretowaną orientacyjnie jako archetypy: gracze defensywni,
uniwersalni, ofensywni w obrębie danej pozycji), a wyniki zwizualizowano w
przestrzeni dwóch pierwszych składowych PCA. Istotność statystyczną różnic
między skupieniami dla wybranych cech (\texttt{total\_points},
\texttt{minutes}, \texttt{xg\_per\_90}, \texttt{xa\_per\_90},
\texttt{xg\_conceded\_per\_90}, \texttt{team\_xp}, \texttt{opp\_xp}) oceniono
testem Kruskala–Wallisa (nieparametrycznym odpowiednikiem jednoczynnikowej
ANOVA, właściwym dla nienormalnych rozkładów cech), z korekcją na
wielokrotne testowanie metodą Bonferroniego. Siłę efektu wyrażono
współczynnikiem $\eta^2$:
\begin{equation}
    \eta^2 = \frac{H - k + 1}{n - k},
\end{equation}
gdzie $H$ jest statystyką testu Kruskala–Wallisa, $k$ liczbą skupień, a $n$
liczbą obserwacji, przy umownej interpretacji wielkości efektu:
$\eta^2 \approx 0{,}01$ -- mały, $0{,}06$ -- średni, $\geq 0{,}14$ -- duży.

%======================================================================
\chapter{Inżynieria cech}\label{ch:features}

\section{Zasada braku wycieku danych}

Kluczowym założeniem projektowym potoku inżynierii cech jest wykorzystanie
w zbiorze treningowym wyłącznie informacji dostępnych przed rozpoczęciem
kolejki, której punktacja jest przedmiotem predykcji. W praktyce oznacza
to, że wszystkie cechy zbudowane na podstawie statystyk meczowych
(\texttt{fpl\_cols}, \texttt{perf\_cols}) są przesuwane w czasie
(\texttt{shift}) o co najmniej jedną kolejkę, natomiast cechy dostępne
\emph{ex ante} (cena, prawdopodobieństwa bukmacherskie, rating Elo,
pozycja, informacja o grze u siebie) mogą być wykorzystywane bez
przesunięcia -- stanowią one zbiór \texttt{pre\_game\_cols}.

\section{Cechy oparte na historii występów zawodnika}

Klasa \texttt{FeatureEngineer} generuje trzy rodziny cech historycznych, w
każdym przypadku grupując dane po kluczu \texttt{(code, season)} i
sortując chronologicznie po \texttt{kickoff\_time}:

\begin{itemize}
    \item \textbf{Wykładnicza średnia ruchoma (EMA)} o rozpiętości $w$:
    \begin{equation}
        \text{EMA}_t(w) = \left(1-\alpha\right)\cdot \text{EMA}_{t-1}(w) + \alpha \cdot x_{t-1}, \quad \alpha = \frac{2}{w+1},
    \end{equation}
    obliczana z przesunięciem \texttt{shift(1)}, tak aby wartość dla
    kolejki $t$ nie wykorzystywała statystyk z kolejki $t$;
    \item \textbf{Cechy opóźnione (lag)} rzędu $l$: $x_{t-l}$, dla
    $l \in \{1,2,3,5\}$;
    \item \textbf{Odchylenie standardowe kroczące} w oknie $w$, również z
    przesunięciem o jedną kolejkę, dla $w \in \{3,5,8\}$, mierzące
    zmienność (konsystencję) formy zawodnika.
\end{itemize}

Metoda \texttt{features\_integration} łączy wszystkie trzy rodziny cech
(dla zadanych list rozmiarów okien) z podzbiorem cech \emph{pre-game} oraz
zmienną celu, uzupełniając ewentualne braki wynikające z niepełnej historii
(początek kariery zawodnika lub sezonu) wartością zero.

\section{Cechy pochodne (\emph{feature crosses})}

Na bazie połączonych źródeł danych zbudowano dodatkowe cechy pochodne,
łączące informację o formie zawodnika z kontekstem meczowym:

\begin{align}
    \text{transfers\_trend} &= \frac{\text{transfers\_balance}}{\text{selected}}, \\[4pt]
    \text{prob\_win\_context} &= \text{prob\_win}, \\[4pt]
    \text{player\_match\_xg\_expected} &= \text{xg\_per\_90} \cdot \text{prob\_win} \cdot \text{prob\_over\_2.5}, \\[4pt]
    \text{player\_match\_xa\_expected} &= \text{xa\_per\_90} \cdot \text{prob\_win} \cdot \text{prob\_over\_2.5}, \\[4pt]
    \text{match\_value\_efficiency} &= \frac{\text{prob\_win} + 0{,}5 \cdot \text{prob\_draw}}{\text{value}}, \\[4pt]
    \text{hype\_vs\_odds\_ratio} &= \frac{\text{transfers\_trend}}{\text{prob\_win} + 0{,}01}.
\end{align}

Wzór~\eqref{eq:elo} został usunięty wraz z rezygnacją z ratingów Elo. Cechy
\texttt{player\_match\_xg/xa\_expected} łączą jakość ofensywną zawodnika z
kontekstem trudności i charakteru meczu (prawdopodobieństwo wygranej i
liczby goli), natomiast \texttt{hype\_vs\_odds\_ratio} zestawia
popularność transferową zawodnika (\emph{sentyment} społeczności FPL) z
obiektywnym prawdopodobieństwem rynkowym, umożliwiając modelowi wychwycenie
rozbieżności między percepcją a oczekiwaną wydajnością.

\section{Redukcja współliniowości}

Po wygenerowaniu pełnego zestawu cech historycznych wykryto i usunięto
cechy silnie ze sobą skorelowane (próg $|\rho| > 0{,}9$ w macierzy korelacji
Pearsona, z pominięciem zmiennej celu), co ogranicza redundancję informacji
i poprawia stabilność estymacji w modelach liniowych oraz interpretowalność
istotności cech w modelach drzewiastych.

%======================================================================
\chapter{Modele uczenia maszynowego}\label{ch:models}

W pracy porównano trzy rodziny modeli o rosnącej złożoności: uogólnione
modele liniowe, modele zespołowe oparte na drzewach decyzyjnych oraz sieci
neuronowe głębokie, w tym architektury dedykowane danym sekwencyjnym.

\section{Uogólnione modele liniowe (GLM)}

Jako model bazowy (\emph{baseline}) zastosowano uogólniony model liniowy
(bibliotekę \texttt{glum}), z regularyzacją typu \emph{elastic net}:
\begin{equation}
    \hat{\beta} = \arg\min_{\beta} \; -\ell(\beta; \mathbf{y}, \mathbf{X}) + \alpha \left[ \rho \|\beta\|_1 + \frac{1-\rho}{2}\|\beta\|_2^2 \right],
\end{equation}
gdzie $\ell$ jest logarytmem wiarygodności rozkładu z rodziny
wykładniczej. Ze względu na charakter zmiennej celu (nieujemna, silnie
skośna, z nadmiarem zer) rozważono trzy rodziny rozkładów: Poissona,
ujemny dwumianowy (\texttt{negative.binomial}) oraz Tweediego, z funkcją
łącząca dobieraną automatycznie (\texttt{link="auto"}). Model GLM pełni w
pracy rolę punktu odniesienia interpretowalnego -- współczynniki $\beta$
pozwalają ocenić kierunek i siłę wpływu poszczególnych cech.

\section{Modele zespołowe oparte na drzewach}

Do rodziny modeli zespołowych włączono:
\begin{itemize}
    \item \textbf{Random Forest} (\texttt{RandomForestRegressor}) jako
    model referencyjny o niskiej wariancji dzięki agregacji \emph{bagging};
    \item \textbf{XGBoost} (\texttt{XGBRegressor}), z funkcją celu
    \texttt{reg:squarederror} i metryką \texttt{rmse}, obsługą kategorii
    poprzez kodowanie \emph{one-hot} realizowane w potoku
    \texttt{FPLDataPipe};
    \item \textbf{CatBoost} (\texttt{CatBoostRegressor}), z funkcją straty
    \texttt{RMSE} oraz natywną obsługą cech kategorycznych
    (\texttt{cat\_features}) bez konieczności kodowania \emph{one-hot} --
    z tego względu zaimplementowano osobną wersję potoku danych
    (\texttt{FPLDataPipe(catboost\_flag=True)}), przekazującą cechy
    kategoryczne jako łańcuchy znaków (\texttt{passthrough}).
\end{itemize}

Strojenie hiperparametrów obu modeli gradient boostingu realizowane jest
metodą przeszukiwania losowego (\texttt{RandomizedSearchCV}) z
walidacją krzyżową typu \texttt{TimeSeriesSplit} (rozdział
\ref{ch:training}), z rozkładami hiperparametrów przedstawionymi w
tabeli~\ref{tab:hparams}.

\begin{table}[h]
\centering
\caption{Przestrzenie przeszukiwania hiperparametrów modeli zespołowych.}
\label{tab:hparams}
\begin{tabular}{lll}
\toprule
Model & Hiperparametr & Rozkład / zakres \\
\midrule
CatBoost & \texttt{iterations} & $\mathcal{U}\{150,\dots,500\}$ \\
         & \texttt{depth} & $\mathcal{U}\{4,\dots,7\}$ \\
         & \texttt{learning\_rate} & $\mathcal{U}(0{,}01;\,0{,}06)$ \\
         & \texttt{l2\_leaf\_reg} & $\mathcal{U}(1{,}0;\,9{,}0)$ \\
         & \texttt{bagging\_temperature} & $\mathcal{U}(0;\,5)$ \\
         & \texttt{random\_strength} & $\mathcal{U}(0{,}5;\,5{,}5)$ \\
\midrule
XGBoost  & \texttt{n\_estimators} & $\mathcal{U}\{150,\dots,500\}$ \\
         & \texttt{max\_depth} & $\mathcal{U}\{3,\dots,7\}$ \\
         & \texttt{subsample} & $\mathcal{U}(0{,}7;\,1{,}0)$ \\
         & \texttt{colsample\_bytree} & $\mathcal{U}(0{,}6;\,0{,}9)$ \\
         & \texttt{reg\_alpha} / \texttt{reg\_lambda} & $\mathcal{U}(0;\,2)$ / $\mathcal{U}(0{,}5;\,3{,}5)$ \\
\bottomrule
\end{tabular}
\end{table}

\section{Wielowarstwowy perceptron (MLP)}

Sieć \texttt{MLP} operuje na cechach tabularycznych (zagregowanych,
niesekwencyjnych) i zawiera dodatkowy, uczony mechanizm \emph{bramkujący}
selekcję cech: wektor parametrów $\mathbf{g}$ jest przepuszczany przez
funkcję sigmoidalną i mnożony element-po-elemencie przez wejście,
\begin{equation}
    \tilde{\mathbf{x}} = \mathbf{x} \odot \sigma(\mathbf{g}),
\end{equation}
co pozwala sieci w trakcie treningu tłumić cechy mało informatywne. Dalsze
warstwy ukryte mają strukturę
\texttt{Linear} $\rightarrow$ \texttt{LayerNorm} $\rightarrow$
\texttt{LeakyReLU} $\rightarrow$ \texttt{Dropout}, z malejącym
prawdopodobieństwem odrzucenia w kolejnych warstwach.

\section{Sieć konwolucyjna 1D (Conv1D)}

Model \texttt{Conv1DRegressor} operuje na sekwencjach $T$ ostatnich
kolejek zawodnika (oknie przesuwnym, opisanym w sekcji
\ref{sec:seqpipe}), traktując wymiar czasu jako wymiar przestrzenny
konwolucji jednowymiarowej. Architektura opiera się na blokach
resztkowych (\emph{residual}):
\begin{equation}
    \mathbf{h}^{(l+1)} = \text{ReLU}\big(\,\text{Block}(\mathbf{h}^{(l)}) + \text{Skip}(\mathbf{h}^{(l)})\,\big),
\end{equation}
gdzie każdy blok składa się z dwóch warstw \texttt{Conv1d} z
normalizacją wsadową (\texttt{BatchNorm1d}) i aktywacją \texttt{GELU}, a
połączenie skrótowe (\texttt{Skip}) jest projekcją $1\times1$, gdy liczba
kanałów wejściowych i wyjściowych bloku różni się. Po ostatnim bloku
zastosowano globalne uśrednianie (\texttt{AdaptiveAvgPool1d}) wzdłuż
wymiaru czasu, a wynik przekazywany jest do głowicy regresyjnej
(\texttt{Linear} $\rightarrow$ \texttt{LayerNorm} $\rightarrow$
\texttt{LeakyReLU} $\rightarrow$ \texttt{Dropout} $\rightarrow$
\texttt{Linear}).

\section{Sieć rekurencyjna LSTM}

Model \texttt{LSTMRegressor} przetwarza tę samą reprezentację sekwencyjną
co model konwolucyjny, poprzez warstwę projekcji wejścia, wielowarstwową
sieć \texttt{LSTM} (z opcjonalną dwukierunkowością) oraz głowicę
regresyjną operującą na ostatnim stanie ukrytym $h_T$ (lub konkatenacji
ostatnich stanów obu kierunków w wariancie dwukierunkowym). Zarówno model
konwolucyjny, jak i rekurencyjny wykorzystują ten sam mechanizm
uczonej bramki selekcji cech, co model MLP.

\section{Potoki danych}\label{sec:seqpipe}

Dla modeli tabularycznych (GLM, modele drzewiaste, MLP) zaimplementowano
klasę \texttt{FPLDataPipe}, opartą na \texttt{sklearn.ColumnTransformer}:
cechy kategoryczne kodowane są metodą \emph{one-hot} (z odrzuceniem
pierwszej kategorii referencyjnej) lub przekazywane bez zmian dla
CatBoost, natomiast cechy numeryczne standaryzowane są metodą
\texttt{StandardScaler}, dopasowywaną wyłącznie na zbiorze treningowym.

Dla modeli sekwencyjnych (Conv1D, LSTM) zaimplementowano
\texttt{SequenceDataPipe} wraz z klasą \texttt{SequenceDataset},
budującą dla każdego zawodnika i sezonu osobno okna przesuwne długości
$T=8$ kolejek: dla kolejki $t$ oknem wejściowym jest podsekwencja cech z
kolejek $[t-T,\,t-1]$, a etykietą -- \texttt{total\_points} z kolejki
$t$. Grupowanie po parze \texttt{(code, season)} zapobiega tworzeniu okien
przekraczających granicę sezonu (transfer zawodnika, awans/spadek), co
mogłoby wprowadzić nierealistyczną ciągłość formy.

%======================================================================
\chapter{Funkcja straty i metryki specyficzne dla FPL}\label{ch:loss}

\section{Motywacja}

Standardowy błąd średniokwadratowy (MSE) traktuje wszystkich zawodników
symetrycznie, podczas gdy z perspektywy użytkownika FPL błędy predykcji
dla zawodników \emph{premium} (drodzy, kluczowi dla wyniku drużyny) są
znacznie bardziej kosztowne niż dla zawodników rezerwowych, a
niedoszacowanie potencjału zawodnika (\emph{underprediction}) prowadzi do
utraconych okazji transferowych w większym stopniu niż analogiczne
przeszacowanie. Ponadto standardowy MSE nie uwzględnia faktu, że
zawodnicy niewystępujący w danej kolejce (\texttt{minutes}$\,=0$) generują
z definicji zerową punktację, co bez odpowiedniej wagi zniekształca
gradient w kierunku trywialnej predykcji zera.

\section{Funkcja straty \texttt{FPLLoss}}

Zaprojektowana funkcja straty modyfikuje bazowy błąd średniokwadratowy
trzema czynnikami wagowymi:

\begin{equation}
    \mathcal{L} = \frac{1}{N}\sum_{i=1}^{N} \big(y_i - \hat{y}_i\big)^2 \cdot w^{\text{cena}}_i \cdot w^{\text{asym}}_i \cdot \big(\text{minutes}_i + 0{,}1\big),
\end{equation}
gdzie:
\begin{itemize}
    \item \textbf{waga cenowa} $w^{\text{cena}}_i$ podbija istotność błędu
    dla zawodników o cenie powyżej progu \texttt{p\_high} za pomocą miękkiej
    (różniczkowalnej) maski sigmoidalnej:
    \begin{equation}
        w^{\text{cena}}_i = 1 + (w_{\text{premium}}-1)\cdot \sigma\!\big(20\cdot(\text{price}_i - p_{\text{high}})\big);
    \end{equation}
    \item \textbf{waga asymetrii} $w^{\text{asym}}_i$ silniej karze
    niedoszacowanie ($y_i > \hat{y}_i$) niż przeszacowanie, poprzez funkcję
    \texttt{softplus} zastosowaną do różnicy błędu:
    \begin{equation}
        w^{\text{asym}}_i = 1 + (c_{\text{under}}-1)\cdot \frac{\text{softplus}_\beta(y_i-\hat{y}_i)}{\text{softplus}_\beta(y_i-\hat{y}_i)+1}, \quad \beta=2;
    \end{equation}
    \item \textbf{waga minut} $(\text{minutes}_i+0{,}1)$ obniża wpływ
    obserwacji zawodników niewystępujących w meczu, jednocześnie unikając
    zerowania gradientu (stała $0{,}1$ zapobiega całkowitemu wyzerowaniu
    wagi).
\end{itemize}

Progi cenowe $p_{\text{low}}$, $p_{\text{high}}$ wyznaczane są empirycznie
jako, odpowiednio, mediana (50. percentyl) i 95,5. percentyl rozkładu ceny
w zbiorze treningowym, po standaryzacji cechy.

\section{Metryki segmentowane}

Analogicznie do funkcji straty, ewaluacja modeli nie ogranicza się do
zagregowanego MSE/MAE, lecz wykorzystuje autorskie metryki
\texttt{CustomMSE} i \texttt{CustomMAE} (zaimplementowane jako klasy
\texttt{torchmetrics.Metric}), raportujące błąd oddzielnie w trzech
segmentach cenowych:
\begin{equation}
    \text{MSE@}k = \frac{1}{|S_k|}\sum_{i \in S_k} (y_i-\hat{y}_i)^2, \qquad k \in \{\text{budget}, \text{mid}, \text{premium}\},
\end{equation}
gdzie $S_{\text{budget}} = \{i: \text{price}_i < p_{\text{low}}\}$,
$S_{\text{mid}} = \{i: p_{\text{low}} \leq \text{price}_i < p_{\text{high}}\}$,
$S_{\text{premium}} = \{i: \text{price}_i \geq p_{\text{high}}\}$.
Metryki te są zgodne implementacyjnie zarówno z modelami tabularycznymi
(indeks ceny w wymiarze cech), jak i sekwencyjnymi (indeks ceny w
ostatnim kroku czasowym okna), dzięki parametrowi \texttt{model\_type}
(\texttt{"mlp"} / \texttt{"lstm"} / \texttt{"conv1d"}).

Dodatkowo raportowany jest skorygowany współczynnik determinacji
$R^2_{\text{adj}}$:
\begin{equation}
    R^2_{\text{adj}} = 1 - (1-R^2)\cdot\frac{n-1}{n-p-1},
\end{equation}
gdzie $n$ jest liczbą obserwacji, a $p$ liczbą cech, co koryguje $R^2$ o
efekt liczby predyktorów i umożliwia rzetelne porównanie modeli o różnej
wymiarowości wejścia.

%======================================================================
\chapter{Trenowanie modeli i strojenie hiperparametrów}\label{ch:training}

\section{Podział danych}

Podział na zbiór treningowy, walidacyjny i testowy wykonano w sposób
chronologiczny, a nie losowy, co jest metodologicznie konieczne przy
danych szeregów czasowych: sezony 2022/23 i 2023/24 stanowią zbiór
treningowy, sezon 2024/25 -- zbiór walidacyjny, a bieżący sezon 2025/26 --
zbiór testowy. Taki podział odwzorowuje rzeczywisty scenariusz użycia
modelu (predykcja na danych z przyszłości względem danych treningowych) i
eliminuje ryzyko wycieku informacji między sezonami.

\section{Standaryzacja zmiennej celu}

Zmienna celu jest dodatkowo standaryzowana metodą \texttt{StandardScaler}
(dopasowaną na zbiorze treningowym) przed treningiem sieci neuronowych, co
stabilizuje proces optymalizacji; wszystkie raportowane metryki są
obliczane po odwróceniu standaryzacji (\texttt{inverse\_transform}), a więc w
oryginalnej skali punktowej.

\section{Klasa \texttt{Trainer}}

Uniwersalna klasa \texttt{Trainer} obsługuje trening wszystkich trzech
architektur sieci neuronowych (MLP, Conv1D, LSTM) w oparciu o wspólny
interfejs, obejmujący:
\begin{itemize}
    \item pętlę treningową i walidacyjną (\texttt{train\_step},
    \texttt{test\_step}) z akumulacją metryk segmentowanych opisanych w
    rozdziale~\ref{ch:loss};
    \item wczesne zatrzymanie (\emph{early stopping}) na podstawie straty
    walidacyjnej, z parametrem cierpliwości (\texttt{patience}) i
    zachowaniem najlepszego stanu wag modelu;
    \item harmonogram szybkości uczenia \texttt{ReduceLROnPlateau},
    redukujący \emph{learning rate} po braku poprawy straty walidacyjnej
    przez zadaną liczbę epok, z okresem chłodzenia (\texttt{cooldown}).
\end{itemize}

\section{Strojenie hiperparametrów modeli klasycznych}

Dla modeli GLM, XGBoost i CatBoost strojenie hiperparametrów realizowane
jest metodą \texttt{RandomizedSearchCV} z walidacją krzyżową
\texttt{TimeSeriesSplit} (zachowującą porządek czasowy w obrębie kroswalidacji,
w przeciwieństwie do standardowego podziału losowego $k$-krotnego),
metryką optymalizacji \texttt{neg\_mean\_squared\_error} oraz $20$
iteracjami losowania z przestrzeni hiperparametrów.

\section{Śledzenie eksperymentów}

Wszystkie eksperymenty (parametry, metryki epokowe, artefakty modeli)
są rejestrowane w narzędziu MLflow za pośrednictwem klasy
\texttt{MlflowTracker}, w tym metryki wsadowe logowane zbiorczo
(\texttt{log\_metrics} per epoka) w celu ograniczenia liczby wywołań
sieciowych do serwera śledzącego, co miało istotne znaczenie przy
uruchamianiu treningów w środowisku obliczeniowym SSPCloud/Onyxia.

%======================================================================
\chapter{Wyniki i dyskusja}\label{ch:results}

\emph{Niniejszy rozdział zostanie uzupełniony po zakończeniu pełnego cyklu
eksperymentów dla wszystkich modeli opisanych w
rozdziale~\ref{ch:models}.} Zalecana struktura rozdziału:

\begin{itemize}
    \item tabela zbiorcza porównująca modele (GLM, Random Forest, XGBoost,
    CatBoost, MLP, Conv1D, LSTM) pod względem MSE, MAE oraz
    $R^2_{\text{adj}}$, raportowanych łącznie oraz w podziale na segmenty
    cenowe (budget/mid/premium) i pozycje zawodników;
    \item krzywe uczenia (strata treningowa i walidacyjna w funkcji epoki)
    dla modeli sieciowych, ilustrujące działanie mechanizmu wczesnego
    zatrzymania;
    \item analiza istotności cech (\emph{feature importance}) dla modeli
    drzewiastych oraz wartości nauczonej bramki selekcji cech
    ($\sigma(\mathbf{g})$) dla modeli sieciowych;
    \item analiza błędów (rezydua predykcji) w podziale na segmenty
    cenowe, weryfikująca skuteczność funkcji straty \texttt{FPLLoss} w
    ograniczaniu niedoszacowania zawodników premium;
    \item studium przypadku: predykcje dla wybranej kolejki bieżącego
    sezonu, zestawione z rzeczywistymi wynikami.
\end{itemize}

\section*{TODO: warstwa stackingowa (meta-model)}
\addcontentsline{toc}{section}{TODO: warstwa stackingowa (meta-model)}

\noindent\fbox{\begin{minipage}{0.97\textwidth}
\small
\textbf{Status: sekcja planowana, do zaimplementowania i uzupełnienia
wynikami.} Poniżej zebrano ustalenia projektowe dotyczące rozszerzenia
architektury pracy o drugi poziom modelowania (\emph{stacked generalization}),
które zostanie zaimplementowane i zewaluowane przed finalnym oddaniem
pracy.
\end{minipage}}

\paragraph{Koncepcja.}
Modele bazowe opisane w rozdziale~\ref{ch:models} (GLM, Random Forest,
XGBoost, CatBoost, MLP, Conv1D, LSTM), potencjalnie trenowane na różnych
wariantach funkcji straty (rozdział~\ref{ch:loss}), generują zróżnicowane,
częściowo nieskorelowane wzorce błędów predykcji. Planuje się dodanie
meta-modelu (docelowo prosty regularyzowany GLM/Ridge lub nieujemna
kombinacja wypukła predykcji), przyjmującego jako cechy wejściowe
predykcje modeli bazowych, a jako target -- rzeczywistą wartość
\texttt{total\_points}.

\paragraph{Krytyczne ryzyko: wyciek danych między poziomami.}
Naiwne podejście (modele bazowe trenowane łącznie do sezonu 2025/26,
meta-model testowany na tym samym sezonie) prowadzi do wycieku danych --
predykcje modeli bazowych na zbiorze testowym byłyby wtedy \emph{in-sample},
sztucznie zaniżając błąd. Przyjęta docelowa struktura trójpoziomowa:

\begin{itemize}
    \item \textbf{trening modeli bazowych}: sezony 2022/23 + 2023/24;
    \item \textbf{trening meta-modelu}: predykcje \emph{out-of-fold}
    zamrożonych modeli bazowych na sezonie 2024/25 jako cechy, z targetem
    \texttt{total\_points} z tego samego sezonu;
    \item \textbf{ewaluacja końcowa całego potoku} (baza + meta): sezon
    2025/26, przy czym modele bazowe pozostają wytrenowane wyłącznie na
    2022/23--2023/24 (bez doučenia na 2024/25 lub 2025/26 przed testem).
\end{itemize}

\paragraph{Rozszerzenie: walk-forward zamiast pojedynczego holdout.}
Ze względu na ograniczoną liczbę sezonów treningowych dla meta-modelu
(ryzyko przeuczenia przy silnie skorelowanych predykcjach modeli
bazowych), rozważane jest zastąpienie pojedynczego holdout (2024/25)
blokowym schematem walk-forward: trening bazy na 2022/23 $\rightarrow$
predykcje OOF na 2023/24; trening bazy na 2022/23+2023/24
$\rightarrow$ predykcje OOF na 2024/25; połączenie obu zbiorów OOF jako
rozszerzony zbiór treningowy meta-modelu, z zachowaniem porządku
czasowego.

\paragraph{Projekt meta-modelu.}
\begin{itemize}
    \item silna regularyzacja (Ridge / Elastic Net w \texttt{glum}) ze
    względu na współliniowość predykcji modeli bazowych;
    \item wariant alternatywny: nieujemna kombinacja wypukła wag
    ($w_j \geq 0$, $\sum_j w_j = 1$) -- mniej elastyczna, ale bardziej
    interpretowalna i odporna na przeuczenie przy małej liczbie
    obserwacji;
    \item rozszerzenie typu \emph{feature-weighted stacking}: dodanie do
    wejścia meta-modelu kontekstu (segment cenowy, pozycja, minuty z
    poprzedniej kolejki), umożliwiające różnicowanie wag modeli bazowych
    względem segmentu (np. większe zaufanie do CatBoost dla obrońców, do
    LSTM dla napastników).
\end{itemize}

\paragraph{Baseline i istotność statystyczna.}
Cały potok (baza + meta) zostanie porównany nie tylko wewnętrznie, ale
także z: (i) oficjalnym \texttt{xP} (\emph{expected points}) dostarczanym
przez FPL, obecnym już w zbiorze danych, jako punktem odniesienia
praktycznej użyteczności, oraz (ii) prostym uśrednieniem predykcji
modeli bazowych, jako punktem odniesienia dla oceny, czy nauczony
meta-model wnosi wartość ponad naiwny blend. Różnice w metrykach błędu
między wariantami zostaną ocenione testem istotności statystycznej
(test Diebolda--Mariano lub bootstrap na poziomie kolejek), ze względu na
ograniczoną liczbę sezonów testowych.

\paragraph{Uwagi dot. wdrożenia (poza zakresem ewaluacji offline).}
Docelowo modele bazowe i meta-model będą wersjonowane łącznie w MLflow
Model Registry (retrening meta-modelu przy każdej zmianie warstwy bazowej),
a predykcje generowane cyklicznie (raz na kolejkę) w ramach zaplanowanego
zadania wsadowego, bez potrzeby wnioskowania w czasie rzeczywistym.


\chapter{Wybór składu --- problem plecakowy}\label{ch:squad}

%======================================================================
\chapter{Decyzje dotyczące składu i system rekomendacji}\label{ch:recommender}
%======================================================================
\chapter*{Podsumowanie}
\addcontentsline{toc}{chapter}{Podsumowanie}

W pracy zaprojektowano i zaimplementowano kompletny potok predykcji
punktacji zawodników Fantasy Premier League, obejmujący automatyczne
pozyskiwanie i integrację danych z trzech niezależnych źródeł,
eksploracyjną analizę danych z wykorzystaniem metod statystycznych
odpornych na nienormalność rozkładów (korelacja Spearmana, test
Kruskala–Wallisa), inżynierię cech respektującą rygor braku wycieku
informacji z przyszłości oraz porównanie trzech rodzin modeli uczenia
maszynowego o rosnącej złożoności. Istotnym wkładem pracy jest
zaprojektowanie funkcji straty \texttt{FPLLoss} oraz towarzyszących jej
metryk segmentowanych względem ceny rynkowej zawodnika, odzwierciedlających
asymetryczną użyteczność decyzyjną błędu predykcji z perspektywy
rzeczywistego użytkownika gry.

Do naturalnych kierunków dalszego rozwoju systemu należą: (i) budowa
modelu zespołowego (\emph{ensemble}) łączącego predykcje modeli
tabularycznych i sekwencyjnych, (ii) uwzględnienie dodatkowych źródeł
danych (np. statystyk zaawansowanych typu \emph{expected threat},
danych o składach wyjściowych ogłaszanych przed meczem), (iii)
przeniesienie modelu do trybu wnioskowania \emph{online}, z cotygodniową
douczaniem na najświeższych danych bieżącego sezonu, oraz (iv)
rozszerzenie zadania predykcyjnego o element optymalizacji doboru składu
drużyny w ramach ograniczeń budżetowych gry (zadanie programowania
całkowitoliczbowego na predykcjach modelu).

 \begin{thebibliography}{9}
 \bibitem{panek}
 E. Panek, \emph{Ekonomia matematyczna}, Akademia Ekonomiczna w Poznaniu,
 Poznań 2003.
 \bibitem{fpl_vaastav}
 Vaastav Anand, \emph{Fantasy-Premier-League} -- repozytorium danych
 historycznych FPL, \url{https://github.com/vaastav/Fantasy-Premier-League}.
 \bibitem{football_data}
 \emph{Football-Data.co.uk} -- historyczne kursy bukmacherskie lig
 piłkarskich, \url{https://www.football-data.co.uk}.
 \bibitem{tobyddb}
 \ldots
 \end{thebibliography}
  
 \end{document}